problem:  
Let \( M^{2n} \) be a compact, simply-connected Riemannian manifold with nonnegative sectional curvature and an effective, isometric \( T^n \)-action of maximal rank. It is known (Grove–Petersen–Wilking) that \( M/T^n \) is an Alexandrov space with curvature \( \ge 0 \). In classical toric topology, each **smooth torus manifold** with orbit space a *simple convex polytope* is determined by a *characteristic function* labeling the facets.  

**Problem:**  
Assume further that the orbit space \( Q = M/T^n \) is homeomorphic to a simple convex polytope and inherits an Alexandrov metric \( d_Q \) from \( M \).  
1. Can the Alexandrov metric \( d_Q \) (together with the combinatorial type of \( Q \)) uniquely determine the equivalence class of characteristic function that defines the smooth structure of \( M \)?  
2. Equivalently, does the intrinsic curvature structure of \( Q \) constrain or even determine the characteristic data that encode the \( T^n \)-manifold topology?  
3. In particular, under what geometric conditions on \( d_Q \) (e.g., curvature distribution, presence of totally geodesic faces) is \( M \) forced to be equivariantly diffeomorphic to a generalized Bott manifold?  

Why is it a "good" problem:  
This problem precisely targets the **interface among Alexandrov geometry, toric topology, and Riemannian geometry**: it asks whether the nonnegative curvature metric on the orbit space encodes the combinatorial and algebraic data governing torus manifolds. The question is genuinely open because current classification theories rely on purely combinatorial characteristic functions, with no understanding of how curvature conditions on the quotient metric constrain or determine them.  

- **Profound insight:** It probes whether Alexandrov geometric features—curvature bounds of orbit spaces—carry enough information to recover the topological or smooth structure of the total space, a deep and natural unification of metric geometry and toric topology.  
- **Bridge between fields:** It integrates Riemannian geometry (nonnegative curvature), transformation group theory (maximal torus symmetry), and polyhedral/Kähler geometry (toric manifolds). Progress would require developing new links between curvature and combinatorial data.  
- **Extreme simplicity:** The statement is short and conceptually clear: *does the quotient metric determine the manifold?*—a visually and geometrically compelling question hiding significant depth.  
- **Unknown and deep:** No known theorem currently relates the Alexandrov metric of \( M/T^n \) to the smooth torus action characteristic data. Answering it would open a new research direction in “metric rigidity of torus manifolds” under curvature constraints.